% =============================================================================
%  KAYNAKLAR / REFERENCES  -- APA 7
% -----------------------------------------------------------------------------
%  Her kaynağın arasına BİR BOŞ SATIR bırakın. Asılı girinti otomatiktir.
%  Leave ONE BLANK LINE between entries. The hanging indent is automatic.
%  Alfabetik sıraya siz karar verin; sınıf sıralama yapmaz.
%  You control the ordering; the class does not sort.
%
%  Dergi adı eğik / italic journal name:  \textit{Dergi Adı, 7}(3), 193.
%  Bağlantı / a link:                     \url{https://doi.org/...}
%  Ve işareti kaçışlıdır / ampersand:     \&
% =============================================================================
\begin{sosareferences}

Alum, E. U., \& Ugwu, O. P.-C. (2025). Artificial intelligence in disease
diagnosis and treatment. \textit{Discover Applied Sciences, 7}(3), 193.
\url{https://doi.org/10.1007/s42452-025-06616-y}

Aras, S., Drakos, C., Manimangalam, V., \& Nasir, M. A. (2025). Public
acceptance of artificial intelligence (AI) in healthcare.
\textit{Frontiers in Digital Health, 7}, 1664345.
\url{https://doi.org/10.3389/fdgth.2025.1664345}

Atalay, H. N., \& Yücel, Ş. (2024). Decoding privacy concerns in health data
disclosure. \textit{Archives of Public Health, 82}(1), 189.
\url{https://doi.org/10.1186/s13690-024-01416-z}

Beets, B., Newman, T. P., Howell, E. L., Bao, L., \& Yang, S. (2023).
Surveying public perceptions of artificial intelligence in health care in the
United States: Systematic review. \textit{Journal of Medical Internet
Research, 25}(1), e40337. \url{https://doi.org/10.2196/40337}

Esmaeilzadeh, P., Mirzaei, T., \& Dharanikota, S. (2021). Patients' perceptions
toward human--artificial intelligence interaction in health care: Experimental
study. \textit{Journal of Medical Internet Research, 23}(11), e25856.
\url{https://doi.org/10.2196/25856}

\end{sosareferences}
